% Flavor catalog per-process template.
% process_id: E008
% family: edm_neutrino
% owner_agent_id: PKA-E008
% writer_agent_id: null
% checker_agent_id: null
% opus_signoff_id: null
% Canonical metadata sidecar: E008.yaml
% Status is controlled by the YAML sidecar status_history, not this comment.

\section{E008: Quark chromo-EDM bounds}

\subsection*{Process}
\textbf{Standard notation:} \(\tilde d_q\), the flavor-diagonal quark
chromo-electric dipole moment in the CP-odd operator
\(\bar q T^a \sigma^{\mu\nu}\gamma_5 q\,G^a_{\mu\nu}\).  This entry covers
EFT-level bounds on light-quark qCEDMs inferred from neutron and
\(^{199}\mathrm{Hg}\) EDM searches, not a directly measured particle property.

\subsection*{PDG / equivalent value(s)}
The experimental anchors are the PDG Live 2026 neutron limit
\[
  |d_n| < 1.8\times10^{-26}\ e\,\mathrm{cm}\quad(90\%~\mathrm{CL})
\]
from Abel 2020, and the Graner et al. 2016 mercury result
\[
  d_{\rm Hg}=(2.20\pm2.75_{\rm stat}\pm1.48_{\rm syst})\times10^{-30}
  \ e\,\mathrm{cm},\qquad
  |d_{\rm Hg}|<7.4\times10^{-30}\ e\,\mathrm{cm}\quad(95\%~\mathrm{CL}).
\]
Using the Pospelov--Ritz QCD-sum-rule relation with the quark EDMs set to
zero gives the central neutron translation
\[
  |\tilde d_d+0.5\,\tilde d_u| < 1.6\times10^{-26}\ \mathrm{cm},
\]
with the quoted \((1\pm0.5)\) hadronic normalization uncertainty.  Using the
standard mercury estimate
\[
  d_{\rm Hg}=7\times10^{-3}e(\tilde d_u-\tilde d_d)+10^{-2}d_e+\cdots
\]
and setting non-qCEDM sources to zero gives the central equivalent bound
\[
  |\tilde d_u-\tilde d_d| < 1.1\times10^{-27}\ \mathrm{cm}.
\]
The local snapshots are
\texttt{flavor\_catalog/references/E008/pdg2026\_neutron\_edm\_datablock.txt}
and \texttt{flavor\_catalog/references/E008/graner2016\_hg199\_arxiv1601\_04339.txt};
the translation formulae are in the Pospelov--Ritz and
Olive--Pospelov--Ritz--Santoso snapshots.

\subsection*{Relevance to the RS / anarchic-flavor pipeline}
Anarchic warped models contain new CP phases and chirality flips in KK-fermion,
Higgs, and gauge loops.  Even when tree-level \(\Delta F=2\) mixing is under
control, these loops can generate flavor-diagonal quark EDMs and qCEDMs.  The
qCEDM lane is therefore a direct probe of RS KK-loop CP violation and
wrong-chirality Yukawa structures.

\subsection*{Post-2008 developments}
Relative to the CFW 2008 RS-flavor baseline, the experimental inputs tightened:
the direct neutron anchor is now the 2020 PDG limit above, and the 2016
\(^{199}\mathrm{Hg}\) result improved the prior mercury limit by a factor of
4.  Koenig, Neubert, and Straub 2014 recast electromagnetic and chromomagnetic
quark dipoles in partial-compositeness and warped-language Wilson coefficients,
finding strong neutron-EDM pressure when wrong-chirality Yukawas are present.
On the theory side, Bhattacharya et al. 2024 target the isovector qCEDM
neutron matrix element with lattice QCD and explicit control of power-divergent
operator mixing.

\subsection*{Constraint validity and model dependence}
The direct neutron and mercury limits are robust null measurements.  The qCEDM
bounds are not standalone PDG averages: they assume a low-energy CP-odd basis,
neglect cancellations with quark EDMs, semileptonic sources, four-quark
operators, \(\bar\theta\), and the Weinberg operator, and inherit QCD, nuclear,
and atomic matrix-element uncertainties.  Mercury is especially powerful for
the isovector combination but is also more nuclear-model dependent.

\subsection*{Code coverage in this repo}
\textbf{NO.}  Required greps over \texttt{quarkConstraints/}, \texttt{qcd/},
\texttt{flavorConstraints/}, \texttt{neutrinos/}, \texttt{yukawa/},
\texttt{warpConfig/}, \texttt{solvers/}, \texttt{scanParams/}, and
\texttt{tests/} found no qCEDM, neutron EDM, mercury EDM, or Weinberg-operator
constraint.  The only nearby dipole implementation is the unrelated
\(\mu\to e\gamma\) helper in \texttt{flavorConstraints/muToEGamma.py:3},
\texttt{:21}, and \texttt{:81}.

\subsection*{Implementation difficulty}
\textbf{HIGH.}  A live E008 constraint would need new CP-odd quark EDM/qCEDM
operators, loop matching from the RS spectrum, QCD running and threshold
mixing, plus hadronic and nuclear matrix-element choices.  The existing
\(\Delta F=2\) SLL/SLR/VLL/VRR/LR basis does not cover this observable.

\subsection*{Key references}
Process-local source keys before bibliography consolidation:
\texttt{PDG2026:NeutronEDM}, \texttt{Graner:HgEDM2016},
\texttt{PospelovRitz:NeutronCEDM2000},
\texttt{OlivePospelovRitzSantoso:HgCEDM2005},
\texttt{CsakiFalkowskiWeiler:RSFlavor2008},
\texttt{KoenigNeubertStraub:DipoleComposite2014}, and
\texttt{Bhattacharya:IsoVectorQCEDM2024}.
